\documentclass[a4paper,french]{article}

\usepackage{cmap}
\usepackage[utf8x]{inputenc}
\usepackage[T1]{fontenc}

\usepackage[francais,bloc]{automultiplechoice}
\usepackage{babel}
\frenchsetup{ListItemsAsPar=true, AutoSpacePunctuation, ThinColonSpace, INGuillSpace, InnerGuillSingle, SuppressWarning}

% à déclarer après babel et automultiplechoice
\usepackage{cleveref}
\usepackage{expl3}

\begin{document}

\setdefaultgroupmode{withoutreplacement}
%%% préparation des groupes

% Insertion des groupes qui peuvent s'écrire directement dans le fichier.
% Pour simplifer le code l'inclusion est privilégiée.
\element{geographie}{
	\begin{question}{Paris}\AMClabel{Paris}
		Dans quel continent se situe Paris~?
		\begin{reponseshoriz}
			\bonne{L'Europe}
			\mauvaise{L'Afrique}
			\mauvaise{L'Asie}
			\mauvaise{La planète Mars}
		\end{reponseshoriz}
	\end{question}
}

\element{geographie}{
	\begin{question}{Cameroun}\AMClabel{Cameroun}
		Quelle est la capitale du Cameroun~?
		\begin{reponseshoriz}
			\bonne{Yaoundé}
			\mauvaise{Douala}
			\mauvaise{Abou-Dabi}
		\end{reponseshoriz}
	\end{question}
}

\element{geographie}{
	\begin{question}{Loire}\AMClabel{Loire}
		Quel est le plus grand fleuve français ?
		\begin{reponseshoriz}
			\bonne{Loire}
			\mauvaise{Rhône}
			\mauvaise{Seine}
			\mauvaise{Rhin}
		\end{reponseshoriz}
	\end{question}
}

\element{geographie}{
	\begin{question}{Rhin}\AMClabel{Rhin}
		Quel est le fleuve français dont la source n'est  pas en France ?
		\begin{reponseshoriz}
			\bonne{Rhin}
			\mauvaise{Rhône}
			\mauvaise{Seine}
			\mauvaise{Loire}
		\end{reponseshoriz}
	\end{question}
}

\element{histoire}{
	\begin{question}{RF}\AMClabel{RF}
		En quelle année débute la Révolution française ?
		\begin{reponseshoriz}
			\bonne{1789}\mauvaise{1492}\mauvaise{1815}\mauvaise{1914}
		\end{reponseshoriz}
		
	\end{question}
}

\element{histoire}{
	\begin{question}{Napoleon}\AMClabel{Napoleon}
		En quelle année Napoléon est-il sacré empereur ?
		\begin{reponseshoriz}
			\bonne{1804}\mauvaise{1812}\mauvaise{1830}\mauvaise{1870}
		\end{reponseshoriz}
		
	\end{question}
}

\element{histoire}{
	\begin{question}{armistice}\AMClabel{armistice}
		En quelle année l'armistice met-il fin à la Première Guerre mondiale ?
		\begin{reponseshoriz}
			\bonne{1918}\mauvaise{1945}\mauvaise{1916}\mauvaise{1939}
		\end{reponseshoriz}
		
	\end{question}
}


%%% fabrication des copies

\exemplaire{5}{

%%% debut de l'en-tête des copies :

\noindent{\bf QCM  \hfill TEST}
\vspace*{.5cm}

\hspace*{\fill}\champnom{\fbox{\begin{minipage}{.5\linewidth}
Nom et prénom :

\vspace*{.5cm}\namefielddots
\vspace*{1mm}
\end{minipage}}}
\medskip

\vspace{2ex}
\noindent \textbf{Note  :}
\begin{itemize}
\item Les \AMCrefs{cref}{Paris,Cameroun,Loire,Rhin} font appel à vos connaissances en géographie.

\item Pour l'histoire : Q. \AMCrefs{labelcref}{RF,Napoleon,armistice}, \AMCrefs{cpageref}{RF,Napoleon,armistice}.
\end{itemize}
\vspace{2ex}

\noindent\hrulefill
\vspace{4ex}

\cleargroup{tout}
\copygroup[3]{geographie}{tout}
\copygroup[2]{histoire}{tout}
\restituegroupe{tout}

}

\end{document}